\documentclass[11pt]{article}
\usepackage[a4paper,margin=25mm]{geometry}
\usepackage{amsmath,amssymb,amsthm,booktabs,microtype}
\usepackage[T1]{fontenc}
\usepackage{lmodern}
\usepackage[dvipsnames]{xcolor}
\usepackage[colorlinks=true,linkcolor=MidnightBlue,citecolor=MidnightBlue,urlcolor=MidnightBlue]{hyperref}
\usepackage{fancyhdr}
\pagestyle{fancy}\fancyhf{}
\fancyhead[L]{\small Berge saturation of the subdivided star}
\fancyhead[R]{\small Dev D. Goyal}
\fancyfoot[C]{\thepage}
\setlength{\headheight}{14pt}
\hypersetup{pdftitle={Exact Berge saturation numbers of the subdivided star in every uniformity and its minimality among trees},
pdfauthor={Dev D. Goyal},
pdfsubject={Fully autonomous agentic research},
pdfkeywords={Berge hypergraphs, saturation numbers, trees, subdivided star, uniform hypergraphs, Lean 4}}
\newtheorem{theorem}{Theorem}
\newtheorem{lemma}[theorem]{Lemma}
\newtheorem{proposition}[theorem]{Proposition}
\newtheorem{corollary}[theorem]{Corollary}
\theoremstyle{definition}
\newtheorem{definition}[theorem]{Definition}
\theoremstyle{remark}
\newtheorem{remark}[theorem]{Remark}
\newcommand{\sat}{\operatorname{sat}}
\newcommand{\BB}{\text{Berge-}B_a}
\title{\textbf{Exact Berge saturation numbers\\of the subdivided star in every uniformity\\and its minimality among trees}}
\author{Dev D. Goyal}
\date{1 November 2026}
\begin{document}
\maketitle
\begin{abstract}
Let $B_a$ be the star $K_{1,a}$ with one edge subdivided. For all
$a,r\ge3$ and $n\ge(r-1)a+1$ we prove
\[
\sat_r(n,\BB)=\left\lfloor\frac{a(n+1)}{(r-1)a+1}\right\rfloor ,
\]
where $\sat_r(n,\text{Berge-}F)$ is the least number of hyperedges of an
$r$-uniform hypergraph on $n$ vertices that contains no Berge copy of $F$
but contains one after the addition of any absent $r$-set. Disjoint loose
stars attain the value; for smaller $n$ the saturation number is $a$ or
$\binom nr$. Saturated hypergraphs within $k$ hyperedges of the minimum
consist, for $n\ge((r-1)a+1)((r-1)(k+1)+1)$, of loose stars with at least
$a$ spokes and at most $(r-1)((r-1)a+1)(k+1)-1$ further vertices; the
constant is sharp for $k=0$. In the three-uniform case, for every $a\ge5$
every tree $T\not\cong B_a$ on $a+2$ vertices satisfies
$\sat_3(n,\text{Berge-}T)\ge\lceil(a+1)(n-2)/(2a+3)\rceil$, so $B_a$ is the
unique tree of its order with the least asymptotic lower saturation
coefficient, $a/(2a+1)$, and the unique pointwise minimizer for
$n\ge3(a+1)(2a+1)$. This is the Berge analogue of the theorem of Faudree,
Faudree, Gould and Jacobson that the subdivided star minimizes graph
saturation among trees. Lean~4 verifies the exact formula, with the
saturation number defined as a minimum over all hypergraphs on $n$ labelled
vertices, and the near-minimum structure theorem; the smaller orders, the
sharpness constructions and the comparison among trees have conventional
proofs.
\end{abstract}

\section{Introduction and results}

All hypergraphs are finite, simple and $r$-uniform with $r\ge3$; the vertex
set is part of the data, so isolated vertices count. A \emph{Berge copy} of
a graph $F$ in a hypergraph $H$ consists of an injection $\phi$ of $V(F)$
into $V(H)$ and an injection $\psi$ of $E(F)$ into $E(H)$ with
$\{\phi(u),\phi(v)\}\subseteq\psi(uv)$ for every edge $uv$. A hypergraph is
\emph{Berge-$F$-saturated} if it contains no Berge copy of $F$ but $H+e$
contains one for every $r$-set $e\notin E(H)$; no condition such as
linearity is imposed on $e$. The saturation number $\sat_r(n,\text{Berge-}F)$
is the least number of hyperedges of a Berge-$F$-saturated hypergraph on $n$
vertices. English, Graber, Kirkpatrick, Methuku and Sullivan introduced the
problem and studied paths, cycles, stars, triangles and
matchings~\cite{EGKMS,EGKMS-J}; English, Gerbner, Methuku and Tait proved
linearity for $3\le r\le5$ and for every $r$ when $F$ has a tree
component~\cite{EGMT}, and a recent preprint of Xie proves linearity for
every graph and every uniformity~\cite{Xie}. Exact values are known for few
families: Berge stars~\cite{AE,EGKMS}, matchings and
triangles~\cite{EGKMS}, and recently three-uniform $K_4$~\cite{CHX,CLW};
for paths~\cite{EGKMS} gives bounds that differ by at most three. Clique asymptotics are due to English, Kritschgau,
Nahvi and Sprangel~\cite{EKNS}, and the saturation spectrum of Berge stars
was studied by Bushaw, English, Heath, Johnston and Rombach~\cite{BEHJR}.

For graphs, Faudree, Faudree, Gould and Jacobson~\cite{FFGJ} showed that
among all trees of order $k\ge5$ the star $K_{1,k-2}$ with one edge
subdivided has the least saturation number, and that it is the only tree
attaining the minimum for all $n$. Write $B_a$ for that tree with $a+2$
vertices: centre $c$, stem $b$, terminal $x$ and leaves $y_1,\dots,y_{a-1}$,
with edges $cb$, $bx$ and $cy_i$. It has $a+1$ edges, and for $a\ge3$ it is
not a star. Throughout put
\[
s=r-1,\qquad A=sa+1 .
\]
A \emph{loose star with $t$ spokes} has hyperedges $\{c\}\cup P_i$
($1\le i\le t$) for pairwise disjoint $s$-sets $P_i$ not containing $c$; it
has $st+1$ vertices.

\begin{theorem}\label{thm:main}
Let $a,r\ge3$ and $n\ge A$, and write $n=Aq+t$ with $0\le t<A$. Then
\[
\sat_r(n,\BB)=\left\lfloor\frac{a(n+1)}{A}\right\rfloor
=aq+\left\lfloor\frac ts\right\rfloor
=\left\lceil\frac{a(n-s+1)}{A}\right\rceil .
\]
The value is attained by $q$ disjoint loose stars, one with
$a+\lfloor t/s\rfloor$ spokes and the others with $a$, together with
$t\bmod s$ isolated vertices.
\end{theorem}

\begin{theorem}\label{thm:small}
Let $a,r\ge3$ and $N_0=\max(a+2,r+1)$. Then $\sat_r(n,\BB)=\binom nr$ for
$n<N_0$ and $\sat_r(n,\BB)=a$ for $N_0\le n<A$.
\end{theorem}

For a saturated hypergraph $H$ let $R(H)$ be the number of vertices outside
the components of $H$ that are loose stars with at least $a$ spokes.

\begin{theorem}\label{thm:stability}
Let $a,r\ge3$, $k\ge0$ and $n\ge A(s(k+1)+1)$. If $H$ is a
Berge-$B_a$-saturated hypergraph on $n$ vertices with
$|E(H)|\le\sat_r(n,\BB)+k$, then $H$ has a component that is a loose star
with at least $a$ spokes, and $R(H)\le sA(k+1)-1$.
\end{theorem}

For $k=0$ the threshold is $n\ge rA$ and the bound is $R(H)\le sA-1$; both
are attained (Proposition~\ref{prop:sharp}).

\begin{theorem}\label{thm:trees}
Let $a\ge5$ and let $T$ be a tree on $a+2$ vertices. For every $n\ge1$,
\[
\sat_3(n,\text{Berge-}T)\ge\left\lceil\frac{a(n-2)}{2a+1}\right\rceil ,
\qquad\text{and if }T\not\cong B_a,\qquad
\sat_3(n,\text{Berge-}T)\ge\left\lceil\frac{(a+1)(n-2)}{2a+3}\right\rceil .
\]
Consequently $\sat_3(n,\BB)/n\to a/(2a+1)$, while
$\liminf_n\sat_3(n,\text{Berge-}T)/n\ge(a+1)/(2a+3)$ for every other tree
$T$ of the same order; the gap is $1/((2a+1)(2a+3))$. Moreover
$\sat_3(n,\text{Berge-}T)>\sat_3(n,\BB)$ for every $T\not\cong B_a$ and
every $n\ge3(a+1)(2a+1)$.
\end{theorem}

For $a=3$ and $a=4$ a direct enumeration of the incidence trees with $a$
triples finds further repeatable blocks, for the paths on five and six
vertices, so the restriction $a\ge5$ is used in Section~\ref{sec:trees}.

Theorems~\ref{thm:main} and~\ref{thm:stability} are verified in Lean~4
(Section~\ref{sec:lean}); Theorems~\ref{thm:small} and~\ref{thm:trees} and
Proposition~\ref{prop:sharp} have the conventional proofs given here.

\section{Components and incidence counting}

The \emph{incidence graph} of $H$ is the bipartite graph on
$V(H)\sqcup E(H)$ joining $v$ to $e$ when $v\in e$. Components of $H$ are
the vertex classes of its components; an isolated vertex is a component
without hyperedges. A connected hypergraph with $v$ vertices and $m\ge1$
hyperedges has an incidence graph with $v+m$ nodes and $rm$ edges, so
\[
v\le sm+1,\qquad\text{with equality exactly when the incidence graph is a
tree.}
\]
Call such a component an \emph{incidence tree}. Single hyperedges and loose
stars are incidence trees. In an incidence tree with $m$ hyperedges the sum
of $d(v)-1$ over its vertices is $rm-v=m-1$, so at most $m-1$ vertices have
degree at least two, and at least $(r-2)m+2$ vertices have degree one. Two
hyperedges of an incidence tree meet in at most one vertex.

Every Berge copy of a connected graph lies in a single component, and the
centre of a Berge copy of $B_a$ has degree at least $a$ in $H$, because the
$a$ hyperedges representing the edges at the centre are distinct and contain
it; its stem has degree at least two.

\begin{lemma}\label{lem:tree}
Let $H$ be Berge-$B_a$-saturated with $a,r\ge3$. Every incidence-tree
component of $H$ with at least two hyperedges is a loose star with at least
$a$ spokes.
\end{lemma}

\begin{proof}
Let $C$ be such a component with $m\ge2$ hyperedges. Suppose some vertex
$c$ of $C$ has degree at least $a$ and shares a hyperedge $e$ with a vertex
$b$ of degree at least two. Let $f\ne e$ contain $b$ and choose
$x\in f\setminus\{b\}$; then $c\notin f$, and $x$ lies in no hyperedge at
$c$, for otherwise the incidence graph would contain a cycle. For the same
reason the other hyperedges at $c$, of which there are at least $a-1$, meet
one another and $e\cup f$ only in $c$, so choosing one vertex other than $c$
from $a-1$ of them gives a Berge copy of $B_a$
with centre $c$, stem $b$ and terminal $x$, which is impossible. Hence if
$C$ has a vertex of degree at least $a$, all its neighbours have degree one
and $C$ is a loose star with at least $a$ spokes. Otherwise every degree in
$C$ is at most $a-1$, and $C$ has at least $(r-2)m+2\ge r$ vertices of
degree one. An $r$-set of them is absent, since every hyperedge of $C$
contains a vertex of degree at least two; adding it raises those degrees to
at most two and leaves every degree below $a$, so no Berge copy of $B_a$
appears, a contradiction.
\end{proof}

For a saturated hypergraph $H$ let $h$ be the number of loose-star
components with at least $a$ spokes, $e_0$ the number of components
consisting of one hyperedge, and $z$ the number of isolated vertices. For
each other component $C$, which contains a cycle in its incidence graph,
put $j(C)=s|E(C)|-|V(C)|\ge0$, and let $S$ be the sum of the $j(C)$.
Summing $s|E(C)|$ over components gives the exact identity
\begin{equation}\label{eq:count}
s\,|E(H)|=n-h-e_0-z+S .
\end{equation}

\begin{lemma}\label{lem:small}
If $H$ is Berge-$B_a$-saturated with $a,r\ge3$, then $e_0\le1$, $z\le s$,
$e_0z=0$, and $z\le s-1$ when $h\ge1$.
\end{lemma}

\begin{proof}
In each case we add an absent $r$-set and find no Berge copy of $B_a$. If
$e_0\ge2$, add an $r$-set meeting two single-hyperedge components: the new
component has maximum degree two, below $a$. The same works for an
isolated vertex and a single-hyperedge component. If $z\ge r=s+1$, add an
$r$-set of isolated vertices, which forms a new single-hyperedge component.
If $h\ge1$ and $z\ge s$, add the centre of a loose star together with $s$
isolated vertices; the result is a loose star with one more spoke, in which
only the centre has degree above one, and a Berge copy of $B_a$ needs two
vertices of degree at least two.
\end{proof}

\section{Constructions}

\begin{lemma}\label{lem:stars}
Let $a,r\ge3$. A disjoint union of loose stars with at least $a$ spokes
each and at most $s-1$ isolated vertices is Berge-$B_a$-saturated.
\end{lemma}

\begin{proof}
Each component has at most one vertex of degree above one, so there is no
Berge copy of $B_a$. Let $e$ be an absent $r$-set. If $e$ contains a
non-centre vertex $b$ of a star with centre $c$, choose $x\in
e\setminus\{b,c\}$. Represent $cb$ by the spoke containing $b$, $bx$ by
$e$, and $a-1$ further edges at $c$ by other spokes, choosing in each a
vertex other than $c$ and $x$; this is possible because each spoke has
$s\ge2$ vertices besides $c$. Otherwise $e$ consists of centres and at most
$s-1$ isolated vertices, so it contains two centres $c,d$. Represent $cd$
by $e$, the pendant edge at $d$ by a spoke at $d$, and the remaining $a-1$
edges at $c$ by spokes at $c$.
\end{proof}

\begin{proof}[Upper bound in Theorem~\ref{thm:main}]
The hypergraph described in the theorem has $Aq+s\lfloor t/s\rfloor+(t\bmod
s)=n$ vertices and $aq+\lfloor t/s\rfloor$ hyperedges, and it is saturated
by Lemma~\ref{lem:stars}. For the identities write $t=sd+\varepsilon$ with
$0\le\varepsilon<s$, so $0\le d\le a$. Then $a(n+1)-A(aq+d)=a(\varepsilon+1)
-d$ lies in $[0,A-1]$, which gives the floor form, and the ceiling form
follows from $a(n-s+1)+A-1=a(n+1)$.
\end{proof}

For Theorem~\ref{thm:small} we use cones. A family $L$ of distinct
$s$-sets is a \emph{hyperforest} if every nonempty subfamily $J$ satisfies
$|\bigcup J|\ge|J|+1$. The \emph{cone} over $L$ with a new root $c$ has the
hyperedges $\{c\}\cup P$ for $P\in L$.

\begin{lemma}\label{lem:rep}
Let $L$ be a hyperforest, $b\in\bigcup L$ and $x\ne b$. There is a member
$P\ni b$ of $L$ such that the members of $L\setminus\{P\}$ have distinct
representatives avoiding $b$ and $x$.
\end{lemma}

\begin{proof}
After $x$ is deleted from every member, each nonempty subfamily still covers
at least as many vertices as it has members, so by Hall's theorem the
members of $L$ have distinct representatives different from $x$. If $b$ is
not a representative, give it to one member containing $b$ instead. The
member $P$ represented by $b$ satisfies the lemma.
\end{proof}

\begin{lemma}\label{lem:cone}
Let $L$ be a hyperforest of $a$ members, all $s$-sets, without isolated
vertices. The cone over $L$ together with at most $s-1$ isolated vertices is
Berge-$B_a$-saturated.
\end{lemma}

\begin{proof}
It has only $a$ hyperedges, so it has no Berge copy of $B_a$. An absent
$r$-set $e$ contains at most $s-1$ isolated vertices and possibly the root,
so it contains a vertex $b$ of $\bigcup L$; choose $x\in e\setminus\{b,c\}$.
Lemma~\ref{lem:rep} gives $P\ni b$ and representatives of the other $a-1$
members avoiding $b,x$. Represent $cb$ by $\{c\}\cup P$, $bx$ by $e$ and the
other edges at $c$ by the corresponding cone hyperedges.
\end{proof}

\begin{lemma}\label{lem:orders}
For every $v$ with $\max(a+1,s+1)\le v\le sa$ there is a hyperforest of $a$
distinct $s$-sets without isolated vertices on exactly $v$ vertices.
\end{lemma}

\begin{proof}
For $v_0=\max(a+1,s+1)$ take $a$ of the $v_0$ cyclic intervals of length $s$
in a cyclically ordered $v_0$-set. They are distinct since $s<v_0$; any
$k\ge1$ of them cover at least $\min(v_0,k+s-1)\ge k+1$ vertices; and they
cover every vertex, since if $v_0=a+1$ only one interval is omitted and
$s\ge2$, and if $v_0=s+1$ any two intervals cover the cycle. While the
family has fewer than $sa$ vertices, some vertex lies in two members;
replace one of its occurrences by a new vertex. The members stay distinct,
no vertex becomes isolated, every union grows by at most one or stays
fixed, so the hyperforest condition persists, and the order grows by one.
\end{proof}

\begin{proof}[Proof of Theorem~\ref{thm:small}]
If $n<a+2$ there is no injection of $V(B_a)$, and if $n\le r$ there is at
most one hyperedge while $B_a$ needs $a+1$; so for $n<N_0$ every
hypergraph is Berge-$B_a$-free and only the complete one is saturated. Let
$N_0\le n<A$. A saturated hypergraph either has an absent $r$-set, whose
addition creates a copy with $a+1$ hyperedges, so it has at least $a$
hyperedges, or it is complete, with $\binom nr\ge n>a$ hyperedges. For the
upper bound, $\max(a+1,s+1)\le n-1\le sa$, so the cone over the hyperforest
of Lemma~\ref{lem:orders} on $n-1$ vertices is saturated with $a$
hyperedges by Lemma~\ref{lem:cone}.
\end{proof}

\section{The lower bound}

\begin{proof}[Lower bound in Theorem~\ref{thm:main}]
Let $H$ be saturated on $n=Aq+t\ge A$ vertices with $m$ hyperedges. By
Lemma~\ref{lem:tree} the loose-star components with at least $a$ spokes
are exactly the incidence-tree components with at least two hyperedges, and
each has at least $A$ vertices, so $h\le q$. If $h\ge1$, then
$e_0+z\le s-1$ by Lemma~\ref{lem:small}, and~\eqref{eq:count} gives
$sm\ge n-q-s+1=saq+t-s+1$, hence
\[
m\ge aq+\left\lceil\frac{t-s+1}{s}\right\rceil=aq+\left\lfloor\frac
ts\right\rfloor .
\]
If $h=0$, then $e_0+z\le s$ and $sm\ge n-s$, so $m\ge aq+\lceil
(q+t)/s\rceil-1\ge aq+\lfloor t/s\rfloor$, because $q+t>s\lfloor t/s\rfloor$
when $q\ge1$.
\end{proof}

\section{Near-minimum structure}

\begin{proof}[Proof of Theorem~\ref{thm:stability}]
Let $m=|E(H)|\le aq+\lfloor t/s\rfloor+k$ and $\varepsilon=t\bmod s$, so
$s\lfloor t/s\rfloor=t-\varepsilon$. By~\eqref{eq:count} with $S\ge0$ and
$e_0+z\le s$,
\[
h\ge n-sm-s\ge q+\varepsilon-s(k+1)\ge1,
\]
since $q\ge s(k+1)+1$. Hence $e_0+z\le s-1$ and $h\ge n-sm-s+1$, while
$Ah\le n-R(H)$. Therefore
\[
R(H)\le n-A(n-sm-s+1)=sAm-san+A(s-1).
\]
Inserting $m\le aq+\lfloor t/s\rfloor+k$ and $n=Aq+t$ gives $R(H)\le
t-A\varepsilon+sAk+A(s-1)\le sA(k+1)-1$, because $t\le A-1$.
\end{proof}

\begin{proposition}\label{prop:sharp}
Let $a,r\ge3$. Let $U$ be the cone over $a$ distinct $s$-sets of which two
share exactly one vertex and which are otherwise pairwise disjoint; it has
$a$ hyperedges and $sa$ vertices.
\begin{enumerate}
\item For $q\ge s$, the union of $q-s+1$ loose stars with $a$ spokes, $s$
copies of $U$ and $s-1$ isolated vertices is a minimum saturated hypergraph
on $Aq+A-1$ vertices with $R=sA-1$.
\item The union of $r$ copies of $U$ and $s$ isolated vertices is a minimum
saturated hypergraph on $rA-1$ vertices with no loose-star component.
\end{enumerate}
\end{proposition}

\begin{proof}
The link of $U$ has one repeated incidence, so it is a hyperforest. In both
hypergraphs every component has at most $a$ hyperedges, so there is no
Berge copy of $B_a$. Let $e$ be absent. If $e$ contains a non-centre vertex
of a loose star, argue as in Lemma~\ref{lem:stars}; if it contains a
non-root vertex $b$ of a copy of $U$ with root $c$, choose $x\in
e\setminus\{b,c\}$ and argue as in Lemma~\ref{lem:cone}, the vertex $x$
possibly lying in another component. Otherwise $e$ consists of centres,
roots and isolated vertices. If it contains two centres or roots $c,d$,
represent $cd$ by $e$, the pendant edge by a hyperedge at $d$, and the
other $a-1$ edges at $c$ by hyperedges at $c$ with distinct representatives
avoiding $d$ and the terminal; these exist since every link is a
hyperforest. In (1) there are at most $s-1$ isolated vertices, so this
covers every case. In (2) the remaining case is one root $c$ and all $s$
isolated vertices. Let $w$ be the shared vertex of the two overlapping
hyperedges $f_1,f_2$ of that copy of $U$, and $x\in f_2\setminus\{c,w\}$.
Represent $cw$ by $f_1$, $wx$ by $f_2$, one further edge at $c$ by $e$, and
the remaining $a-2$ edges at $c$ by the other $a-2$ hyperedges of the copy,
which meet $f_1\cup f_2$ only in $c$.

The orders are $(q-s+1)A+s^2a+s-1=Aq+A-1$ and $rsa+s=rA-1$, and the sizes
are $(q+1)a$ and $ra$, which equal $\sat_r(n,\BB)$ by
Theorem~\ref{thm:main} because $t=A-1=sa$ in both cases. The vertices
outside loose stars number $s\cdot sa+s-1=sA-1$ in (1).
\end{proof}

\section{Minimality among trees}\label{sec:trees}

In this section $r=3$, $a\ge5$ and $T$ is a tree with $a+1$ edges. A Berge
copy of $T$ selects from each hyperedge it uses one pair, the image of an
edge of $T$; conversely, distinct hyperedges with selected pairs forming a
copy of $T$ give a Berge copy. Core degrees are at most host degrees. If a
host has exactly $a+1$ hyperedges, a Berge copy of $T$ uses all of them.
Three-uniform incidence trees with $m$ hyperedges have $2m+1$ vertices, and
connected components that are not incidence trees satisfy $v\le2m$. For an
ordinary tree $R$ with at least two vertices, the handshake identity gives
\begin{equation}\label{eq:leaves}
L(R)=2+\sum_{x\in V(R)}(d_R(x)-2)_+ ,
\end{equation}
where $L(R)$ is the number of leaves.

\subsection{\texorpdfstring{Internally saturated incidence trees with $a$ triples}{Internally saturated incidence trees with a triples}}

Let $H$ be an incidence tree with exactly $a$ triples such that $H+e$
contains a Berge copy of $T$ for every absent triple $e\subseteq V(H)$.
($H$ itself has too few hyperedges to contain one.)

\begin{lemma}\label{lem:three}
No triple of $H$ has three vertices of degree at least two.
\end{lemma}

\begin{proof}
Suppose $h$ does. Removing $h$ from the incidence tree leaves three branches
meeting $h$ in distinct vertices, each with at least one triple; the $a-1\ge4$
other triples lie in these branches, so one branch has $m\ge2$ triples and
$2m+1$ vertices, hence an absent triple $e$ inside it. A Berge copy of $T$
in $H+e$ uses every triple. Every triple other than $h$, including $e$, lies
inside one branch, and each branch contains a used triple whose selected
pair lies in that branch; the selected pair of $h$ meets at most two
branches, so the selected pairs do not form a connected graph.
\end{proof}

Discarding from every triple one vertex of degree one, keeping both
vertices of degree at least two when there are two, turns $H$ into an
ordinary graph $G$ with $a$ edges and $a+1$ vertices. It is connected and
acyclic, so $G$ is a tree and $H$ is its \emph{expansion}, obtained by adding
a private vertex to each edge. The vertices of degree at least two in $H$
are the nonleaf vertices of $G$, with the same degrees.

\begin{lemma}\label{lem:expansion}
$G$ is a star or a double star, and $L(T)\le L(G)$.
\end{lemma}

\begin{proof}
Let $G$ have $L$ leaves and $I=a+1-L$ nonleaf vertices. $H$ has $a+L\ge3$
vertices of degree one, and no triple of $H$ consists of such vertices,
since every edge of the tree $G$ with $a\ge2$ edges has a nonleaf endpoint. Add
a triple $e$ of them. In $H+e$ they have degree at most two and the other
degrees are unchanged, so by~\eqref{eq:leaves} applied to $T$ and the
monotonicity of $(d-2)_+$, $L(T)\le2+\sum_v(d_G(v)-2)_+=L$. If $I\ge3$, add
a triple $f$ of nonleaf vertices of $G$, absent by Lemma~\ref{lem:three}.
In $H+f$ only the $I$ nonleaf vertices have degree at least two, so a
Berge copy of $T$ has at least $a+2-I=L+1$ leaves, a contradiction. Thus
$I\le2$.
\end{proof}

There is one star and there are $\lfloor(a-1)/2\rfloor$ double stars with
$a$ edges, so at most $1+\lfloor(a-1)/2\rfloor=\lceil a/2\rceil$
isomorphism types of $H$ are possible.

\subsection{Three disjoint copies}

\begin{lemma}\label{lem:threecopies}
If the disjoint union $3H$ of three copies of $H$ is Berge-$T$-saturated,
then $T\cong B_a$ and $H$ is a loose star.
\end{lemma}

\begin{proof}
If a triple meeting three different copies is added, a new copy of $T$ uses
it, and its selected pair joins two of the copies; deleting that edge of
$T$ leaves two rooted subtrees, one in each of these copies, rooted at the
endpoints of the pair, and nothing from the third copy.

Suppose first that $H$ is a loose star with centre $c$ and triples
$\{c,x_i,y_i\}$. A connected Berge subtree containing $c$ selects pairs
through $c$, so it is a star centred at $c$. Adding the triple of the three
centres shows that $T$ is two stars with joined centres. $T$ is not the
star $K_{1,a+1}$: adding a triple of three vertices of degree one leaves
the maximum degree $a$. So $T$ is a double star. Now add $\{c,x_1,x_2\}$:
only $c$ has degree at least three, and $x_1,x_2$ have degree two. The two
centres of $T$ have degrees summing to $a+2\ge7$, so one has degree at least
three and is mapped to $c$, and the other has degree two. Hence $T\cong
B_a$.

Suppose next that $H$ is the expansion of a double star with hubs $u,v$,
with $p\ge1$ pendant triples at $u$ and $q\ge1$ at $v$, $p+q=a-1$, and
central triple $\{u,v,w\}$. A nontrivial Berge subtree rooted at $w$ uses
the central triple with pair $wu$ or $wv$ and then only pendant triples at
that hub, each through the hub; it is a path $w$--hub followed by
$\sigma\le\max(p,q)\le a-2$ leaves at the hub. Add the triple of the three
copies of $w$. If one rooted half were a single vertex, the other would
contain all $a$ old triples used, but it contains at most $a-1$; so both
halves are nontrivial, with $\sigma+\tau=a-2$ leaves. If $\sigma,\tau>0$,
$T$ contains a path with five edges. Adding the absent triple $\{u,v,x\}$,
with $x$ in a pendant triple at $u$, leaves only three vertices of degree at
least two, while such a path has four internal vertices; so this is
impossible. Hence $T$ has a vertex of degree $a-1$ with $a-2$ leaves and a
pendant path of length three, and $a-2\le\max(p,q)$ forces $\min(p,q)=1$;
say $p=a-2$, $q=1$, and let $f$ be the pendant triple at $v$. Add a triple
$e$ of three non-hub vertices of pendant triples at $u$ (there are $2p\ge6$
of them). In $H+e$ only $u$ has degree at least three, namely $a-1$, so the
vertex of degree $a-1$ of $T$ is mapped to $u$ and its edges use all $a-1$
triples containing $u$. The two remaining edges of $T$ are consecutive on
the pendant path but must be represented by the disjoint triples $e$ and
$f$, a contradiction.
\end{proof}

\subsection{Proof of Theorem~\ref{thm:trees}}

Let $K$ be Berge-$T$-saturated on $n$ vertices with $m$ hyperedges. A
component with an absent triple inside it has at least $a$ hyperedges, since
adding the triple creates a connected copy with $a+1$ hyperedges within that
component. An incidence-tree component with $b\ge2$ triples has
$\binom{2b+1}3>b$ triples available, so $b\ge a$. Adding a triple on three
isolated vertices, on two vertices of one single-triple component and one of
another, or on an isolated vertex and two vertices of a single-triple
component creates components with at most three hyperedges and changes
nothing else, so $K$ has at most two isolated vertices, at most one
single-triple component, and not both.

Put $\alpha=a/(2a+1)$ and compare each component with density $\alpha$.
Incidence-tree components with $b\ge a$ triples have $b/(2b+1)\ge\alpha$;
other components have at least half as many hyperedges as vertices; two
isolated vertices fall short by $2\alpha$ and a single-triple component by
$3\alpha-1<2\alpha$. Hence $m\ge\alpha(n-2)$, the first inequality.

Now let $T\not\cong B_a$, let $h$ be the number of incidence-tree
components with exactly $a$ triples, and put $\beta=(a+1)/(2a+3)$. Each such
component is internally saturated, so it is the expansion of a star or a
double star, of at most $\lceil a/2\rceil$ types. No type occurs three
times: three isomorphic components form a Berge-$T$-free hypergraph in which
every absent triple creates a copy (the copy lies in the union of the
three components and the new triple), contradicting
Lemma~\ref{lem:threecopies}. So $h\le2\lceil a/2\rceil\le a+1$.

If $h\ge1$, $K$ has at most one isolated vertex. Indeed, let $H$ be the
expansion of $G$ as above, $c$ a nonleaf vertex of $G$ (one exists since
$a\ge2$) and $z_1,z_2$
isolated vertices, and add $\{c,z_1,z_2\}$. A new copy of $T$ uses the $a+1$
hyperedges of the augmented component, and its new pair is $cz_1$ or $cz_2$,
as $z_1z_2$ would be disconnected from the rest. Removing that leaf leaves a
connected selection using every triple of $H$ once. Such a selection keeps
both nonleaf vertices of a central triple, since otherwise one side of the
corresponding edge of $G$ would be cut off, and the hub of every pendant
triple; so it is isomorphic to $G$, and $T$ is $G$ with a leaf added at a
nonleaf vertex, giving $L(T)=L(G)+1$, against Lemma~\ref{lem:expansion}.

Compare with density $\beta$. Incidence-tree components with at least $a+1$
triples and the components that are not incidence trees have density at
least $\beta$; each of the $h$ components with $a$ triples falls short by
$\beta(2a+1)-a=1/(2a+3)$. If $h=0$ the small components fall short by at
most $2\beta$; if $h\ge1$ by at most $\beta$, since a single-triple
component falls short by $3\beta-1=a/(2a+3)<\beta$. So the total shortfall
is at most $2\beta$ if $h=0$, and at most $(h+a+1)/(2a+3)\le2(a+1)/(2a+3)
=2\beta$ if $h\ge1$; in both cases $m\ge\beta(n-2)$.

The limit for $B_a$ follows from this lower bound and the value
$\lfloor a(n+1)/(2a+1)\rfloor$ of Theorem~\ref{thm:main} with $r=3$, and
$\beta-\alpha=1/((2a+1)(2a+3))$. Finally
\[
\frac{(a+1)(n-2)}{2a+3}-\frac{a(n+1)}{2a+1}=\frac{n-(6a^2+9a+2)}{(2a+1)(2a+3)}>0
\]
for $n\ge6a^2+9a+3=3(a+1)(2a+1)$, so the ceiling on the left exceeds
$\lfloor a(n+1)/(2a+1)\rfloor=\sat_3(n,\BB)$. \qed

\section{Formal verification}\label{sec:lean}

The Lean project in \texttt{lean/} pins Lean~4.35.0 and Mathlib revision
\texttt{56a0058}. Hypergraphs are \texttt{UniformHypergraph V r}: a finite set
of finite sets of vertices of cardinality $r$, on an ambient vertex type, so
multiple hyperedges are excluded and isolated vertices are kept.
\texttt{BroomCopy a E} consists of an injection
$\mathtt{Fin}(a+2)\hookrightarrow V$ and an injection of $\mathtt{Fin}(a+1)$
into the hyperedges of $E$, with both endpoints of the $i$th edge of $B_a$ in
the $i$th hyperedge; the edges are $0$--$1$, $1$--$2$ and $0$--$(i+1)$ for
$i\ge2$, which is $B_a$ with centre $0$, stem $1$ and terminal $2$.
\texttt{Saturated a H} requires that there is no such copy and that one
exists after inserting any $r$-set that is not a hyperedge. The saturation
number is
\[
\texttt{satNumber a r n}=\inf\{m:\exists H\ \text{on }\mathtt{Fin}\,n,\
\texttt{Saturated a H}\wedge|E(H)|=m\}.
\]
The declarations \texttt{Berge.satNumber\_eq} and
\texttt{Berge.exact\_statement\_expanded} state Theorem~\ref{thm:main} for
all $a,r\ge3$ and $n\ge(r-1)a+1$, the second with all definitions unfolded,
and \texttt{Berge.near\_minimum\_structure} states
Theorem~\ref{thm:stability} with the remainder defined as the total order of
the components that are not loose stars with at least $a$ spokes. The proofs
follow Sections~2--5: the construction is built on an explicit vertex type
and transported to $\mathtt{Fin}\,n$, and the lower bound formalizes
Lemmas~\ref{lem:tree} and~\ref{lem:small} and identity~\eqref{eq:count}
through the connected components of the incidence graph. A module-level
audit enumerates every declaration of the fourteen project modules and
checks that each depends only on \texttt{propext}, \texttt{Classical.choice}
and \texttt{Quot.sound}; the sources contain no \texttt{sorry} and no
\texttt{native\_decide}. Theorem~\ref{thm:small},
Proposition~\ref{prop:sharp} and Theorem~\ref{thm:trees} are not formalized.

\section{Prior work and limitations}

The definitions are those of~\cite{EGKMS}. Theorem~\ref{thm:trees} is the
three-uniform Berge analogue of Theorem~5 of~\cite{FFGJ}, with the same
minimizer, and it uses the same strategy of comparing tree components; the
ordinary-graph theorem does not imply it, since an added triple can meet
three components and a minimum incidence tree need not be the expansion of
an ordinary tree before Lemma~\ref{lem:expansion} is proved. Linearity of
$\sat_r(n,\text{Berge-}T)$ for trees is due to~\cite{EGMT}. A search of the
Berge-saturation literature through October 2026, including the papers
cited here, found no exact value for $B_a$ and no comparison of Berge
saturation numbers across trees, but the search is not exhaustive and
priority is not claimed. Theorem~\ref{thm:trees} is proved only for $r=3$
and $a\ge5$; we do not know whether $B_a$ minimizes among trees in higher
uniformities, nor the exact values for other trees.

\section*{Research process and AI disclosure}

This is fully autonomous agentic research. Dev D. Goyal is responsible
for publication. No expert validation or external peer review is
claimed, and the novelty assessment is provisional.

The results, their proofs, the original Lean development (for Lean~4.19.0)
and the literature search were produced by autonomous AI research agents in
an earlier session. Its records describe a fresh internal review, which
rederived the comparison among trees and reran the Lean build, as an Astra
review; they do not otherwise name the system used.

Opus 5.5 by Anthropic used with Claude Code assessed the package for this
release. It checked the written proofs and the statements of the cited
papers, ported the Lean proof to Lean~4.35.0, added the definition of the
saturation number as a minimum and the headline statements, ran the fresh
build and axiom audit, reran the upload's finite checks, and wrote this
note. AI checks are not human peer review.

\begin{thebibliography}{99}
\bibitem{AE}
B.~Austhof and S.~English, \emph{Nearly-regular hypergraphs and saturation
of Berge stars}, Electronic Journal of Combinatorics \textbf{26}(4) (2019),
P4.49. \href{https://arxiv.org/abs/1812.00472}{arXiv:1812.00472}.

\bibitem{BEHJR}
N.~Bushaw, S.~English, E.~Heath, D.~P.~Johnston and P.~Rombach,
\emph{The saturation spectrum of Berge stars},
\href{https://arxiv.org/abs/2502.17686}{arXiv:2502.17686} (2025).

\bibitem{CLW}
X.~Cao, Y.~Lai and C.~Wang, \emph{Saturation in three-uniform Berge
hypergraphs: the five-star spectrum and the $K_4$ saturation number},
\href{https://arxiv.org/abs/2610.32353}{arXiv:2610.32353} (2026).

\bibitem{CHX}
Y.~Chen, J.~He and T.~Xie, \emph{Saturation numbers for $3$-uniform
Berge-$K_4$}, \href{https://arxiv.org/abs/2601.18455}{arXiv:2601.18455}
(2026).

\bibitem{EGKMS}
S.~English, N.~Graber, P.~Kirkpatrick, A.~Methuku and E.~C.~Sullivan,
\emph{Saturation of Berge hypergraphs},
\href{https://arxiv.org/abs/1710.03735}{arXiv:1710.03735} (2017).

\bibitem{EGKMS-J}
S.~English, P.~Gordon, N.~Graber, A.~Methuku and E.~C.~Sullivan,
\emph{Saturation of Berge hypergraphs}, Discrete Mathematics
\textbf{342} (2019), 1738--1761.
\href{https://doi.org/10.1016/j.disc.2019.01.031}{doi:10.1016/j.disc.2019.01.031}.

\bibitem{EGMT}
S.~English, D.~Gerbner, A.~Methuku and M.~Tait,
\emph{Linearity of saturation for Berge hypergraphs}, European Journal of
Combinatorics \textbf{78} (2019), 205--213.
\href{https://arxiv.org/abs/1807.06947}{arXiv:1807.06947}.

\bibitem{EKNS}
S.~English, J.~Kritschgau, M.~Nahvi and E.~Sprangel, \emph{Saturation
numbers for Berge cliques}, European Journal of Combinatorics \textbf{118}
(2024), 103911. \href{https://arxiv.org/abs/2301.02973}{arXiv:2301.02973}.

\bibitem{FFGJ}
J.~Faudree, R.~J.~Faudree, R.~J.~Gould and M.~S.~Jacobson,
\emph{Saturation numbers for trees}, Electronic Journal of Combinatorics
\textbf{16} (2009), R91.
\href{https://doi.org/10.37236/180}{doi:10.37236/180}.

\bibitem{Xie}
T.~Xie, \emph{A linear upper bound for Berge saturation numbers},
\href{https://arxiv.org/abs/2610.11409}{arXiv:2610.11409} (2026).
\end{thebibliography}
\end{document}
